\documentclass[10pt,aspectratio=169]{beamer}
\usetheme{default}
\usepackage[T1]{fontenc}
\usepackage{lmodern}
\usepackage{booktabs}
\usepackage{array}
\usepackage{amsmath}

\definecolor{Green}{HTML}{147A5B}
\definecolor{Navy}{HTML}{17324D}
\definecolor{Gold}{HTML}{D49B2A}
\definecolor{Light}{HTML}{F1F4F6}
\definecolor{Gray}{HTML}{667788}
\definecolor{Red}{HTML}{B54242}

\setbeamercolor{normal text}{fg=Navy,bg=white}
\setbeamercolor{structure}{fg=Green}
\setbeamercolor{frametitle}{fg=white,bg=Navy}
\setbeamercolor{block title}{fg=white,bg=Green}
\setbeamercolor{block body}{fg=Navy,bg=Light}
\setbeamercolor{alerted text}{fg=Gold}
\setbeamertemplate{navigation symbols}{}
\setbeamertemplate{itemize item}{\color{Green}$\blacktriangleright$}
\setbeamertemplate{footline}{\hfill\color{Gray}\insertframenumber/\inserttotalframenumber\hspace{0.45cm}\vspace{0.22cm}}
\setlength{\tabcolsep}{4pt}
\renewcommand{\arraystretch}{1.18}
\newcolumntype{L}[1]{>{\raggedright\arraybackslash}p{#1}}

\newcommand{\good}[1]{\textcolor{Green}{\textbf{#1}}}
\newcommand{\bad}[1]{\textcolor{Red}{\textbf{#1}}}
\newcommand{\smallnote}[1]{{\scriptsize\color{Gray}#1}}

\title{SLURRY\_FREE\_ACID Modeling}
\subtitle{Project advancement and current evidence}
\author{Internship project}
\date{July 2026}

\begin{document}

\begin{frame}[plain]
  \begin{beamercolorbox}[wd=\paperwidth,ht=\paperheight,dp=0pt]{frametitle}
    \vspace{2.2cm}\hspace{1.3cm}
    \begin{minipage}{0.76\textwidth}
      {\color{white}\LARGE\bfseries\inserttitle\par}
      \vspace{0.3cm}
      {\color{white!85}\Large\insertsubtitle\par}
      \vspace{0.8cm}
      {\color{white!70}\insertauthor\hfill\insertdate}
    \end{minipage}
    \vfill
    \color{Gold}\rule{\paperwidth}{4pt}
  \end{beamercolorbox}
\end{frame}

\begin{frame}{1. The project now covers three operational cases}
  \begin{columns}[T,onlytextwidth]
    \column{0.31\textwidth}
      \textbf{1. Direct future level}
      \[X_{\leq t}\longrightarrow \hat y_{t+h}\]
      Process measurements predict the future acid level directly.
      \vspace{0.2cm}

      \smallnote{Useful when process measurements are available; evaluated against persistence when $y_t$ is observed.}
    \column{0.31\textwidth}
      \textbf{2. Target-anchored forecast}
      \[X_{\leq t},y_{\leq t}\longrightarrow \widehat{\Delta y}_{t,h}\]
      The model predicts the change, then reconstructs
      $\hat y_{t+h}=y_t+\widehat{\Delta y}_{t,h}$.
      \vspace{0.2cm}

      \smallnote{Best forecasting case, but it requires a current free-acid measurement.}
    \column{0.31\textwidth}
      \textbf{3. Target-free virtual sensor}
      \[X_{\leq t}\longrightarrow \hat y_t\]
      Process measurements estimate the current acid level without any target column.
      \vspace{0.2cm}

      \smallnote{Designed for a new CSV containing only process variables.}
  \end{columns}

  \vfill
  \begin{block}{Important distinction}
    These are different deployment contracts. The target-anchored model cannot be used when the uploaded file contains no SLURRY\_FREE\_ACID observations.
  \end{block}
\end{frame}

\begin{frame}{2. Preparation was designed as a controlled experiment}
  \begin{columns}[T,onlytextwidth]
    \column{0.48\textwidth}
      \textbf{Data and targets}
      \begin{itemize}
        \item 44,640 chronological one-minute records
        \item 21 process measurements
        \item Forecast horizons $h\in\{1,5,10,15\}$ minutes
        \item $y_{t+h}=\mathrm{SLURRY\_FREE\_ACID}_{t+h}$
      \end{itemize}
    \column{0.48\textwidth}
      \textbf{Leakage controls}
      \begin{itemize}
        \item Chronological train/validation/test splits
        \item Lag selection uses training targets only
        \item Imputers, scalers and filters fitted on training only
        \item Structural warm-up removed before splitting
      \end{itemize}
  \end{columns}

  \vspace{0.25cm}
  \begin{block}{Scientific purpose}
    Compare feature hypotheses and horizons fairly, rather than create one oversized feature matrix.
  \end{block}
\end{frame}

\begin{frame}{3. Four feature sets test four hypotheses}
  \small
  \begin{tabular}{>{\bfseries}L{0.15\textwidth}L{0.35\textwidth}L{0.42\textwidth}}
    \toprule
    Set & Contents & Hypothesis \\
    \midrule
    A --- Raw & Current process measurements & The current process state may be sufficient. \\
    B --- Temporal & Raw values plus selected lags, rolling behavior, differences and trends & Process history may reveal delayed responses and recent dynamics. \\
    C --- Process-aware & Raw values plus balances, ratios, thermal gradients and stability indicators & Interpretable process relationships may improve prediction and generalization. \\
    D --- Full & Accepted temporal and process-aware features & Historical and process-relationship evidence may be complementary. \\
    \bottomrule
  \end{tabular}

  \vfill
  \smallnote{Examples include TOTAL\_ACID\_FLOW, ACID\_FLOW\_IMBALANCE and TOTAL\_SLURRY\_M3\_FLOW. Ratios and proxies remain provisional when units or stream correspondence are unconfirmed.}
\end{frame}

\begin{frame}{4. A small model set covered complementary behaviors}
  \small
  \begin{tabular}{>{\bfseries}L{0.18\textwidth}L{0.30\textwidth}L{0.44\textwidth}}
    \toprule
    Model & Quick definition & Why it was included \\
    \midrule
    Ridge & Regularized linear regression & Stable, fast and interpretable baseline for many correlated features. \\
    PLS & Latent-variable linear regression & Classical industrial baseline for strongly collinear process sensors. \\
    Extra Trees & Average of randomized decision trees & Captures nonlinear thresholds and feature interactions. \\
    HistGradient & Sequential histogram-based trees & Compact nonlinear boosting with efficient training. \\
    XGBoost / LightGBM & Optimized boosted-tree ensembles & Strong tabular learners for complex nonlinear relationships. \\
    KNN & Prediction from similar historical process states & Tests a local, just-in-time modeling hypothesis. \\
    Weighted blend & Validation-selected average of two models & Combines complementary linear and nonlinear errors. \\
    \bottomrule
  \end{tabular}

  \vfill
  \smallnote{Search spaces were deliberately small. Selection used chronological validation; the untouched test period was evaluated once.}
\end{frame}

\begin{frame}{5. Direct level models did not beat persistence}
  \begin{columns}[T,onlytextwidth]
    \column{0.58\textwidth}
      \small
      \begin{tabular}{rrrr}
        \toprule
        Horizon & MAE & RMSE & vs. persistence \\
        \midrule
        1 min  & 0.269 & 0.332 & \bad{-67.9\%} \\
        5 min  & 0.275 & 0.338 & \bad{-35.1\%} \\
        10 min & 0.280 & 0.344 & \bad{-22.6\%} \\
        15 min & 0.304 & 0.378 & \bad{-27.5\%} \\
        \bottomrule
      \end{tabular}
      \vspace{0.2cm}

      \smallnote{Validation selected Ridge with Feature Set D at every horizon. Negative improvement means the model is worse than carrying the current target forward.}
    \column{0.38\textwidth}
      \textbf{Interpretation}
      \begin{itemize}
        \item The model learned the operating level.
        \item Persistence already has a powerful anchor: the observed $y_t$.
        \item Direct level learning spent capacity reproducing that level instead of learning movement.
      \end{itemize}
  \end{columns}

  \vfill
  \begin{block}{Decision}
    Reframe the anchored task around the future change, not the absolute level.
  \end{block}
\end{frame}

\begin{frame}{6. Predicting change made the anchored forecast useful}
  \begin{columns}[T,onlytextwidth]
    \column{0.58\textwidth}
      \small
      \begin{tabular}{rrrrr}
        \toprule
        Horizon & MAE & RMSE & $R^2$ & vs. persistence \\
        \midrule
        1 min  & 0.102 & 0.151 & 0.889 & \good{+23.2\%} \\
        5 min  & 0.113 & 0.169 & 0.862 & \good{+32.3\%} \\
        10 min & 0.126 & 0.182 & 0.839 & \good{+34.9\%} \\
        15 min & 0.176 & 0.237 & 0.726 & \good{+19.8\%} \\
        \bottomrule
      \end{tabular}
      \vspace{0.2cm}

      \smallnote{Selected model: Ridge, hybrid full set, 188 inputs, $\alpha=10$. Metrics are computed after reconstructing $\hat y_{t+h}=y_t+\widehat{\Delta y}_{t,h}$.}
    \column{0.38\textwidth}
      \textbf{Why it improved}
      \begin{itemize}
        \item $y_t$ supplies the current level.
        \item The model focuses on the smaller future movement.
        \item Target history and process features both added value.
        \item Ten minutes gave the largest relative gain, but with higher absolute error than 1 or 5 minutes.
      \end{itemize}
  \end{columns}

  \vfill
  \alert{Constraint:} this model requires current and past SLURRY\_FREE\_ACID values at prediction time.
\end{frame}

\begin{frame}{7. Process information and target history were complementary}
  \begin{columns}[T,onlytextwidth]
    \column{0.53\textwidth}
      \textbf{Ablation findings}
      \begin{itemize}
        \item Accepted features improved over raw-only inputs by \textbf{10.3--24.2\%}, depending on horizon.
        \item Adding target history to exogenous full features improved RMSE by \textbf{8.2--12.7\%}.
        \item Adding process/full information over target history alone improved RMSE by \textbf{5.4--23.5\%}.
      \end{itemize}
    \column{0.43\textwidth}
      \textbf{Alternatives tested}
      \begin{itemize}
        \item Extra Trees did not replace Ridge after validation.
        \item Transition weighting slightly helped transition errors but worsened overall RMSE.
        \item A logistic warning classifier had low precision (0.14--0.16) and high false-warning rates (0.26--0.29).
      \end{itemize}
  \end{columns}

  \vfill
  \begin{block}{Meaning}
    The compact hybrid model was retained. The warning classifier remains a research diagnostic, not an operational alarm.
  \end{block}
\end{frame}

\begin{frame}{8. A target-free virtual sensor works, with limitations}
  \begin{columns}[T,onlytextwidth]
    \column{0.55\textwidth}
      \textbf{Selected model}
      \begin{itemize}
        \item 50/50 validation-selected blend
        \item Feature Set D Ridge + LightGBM
        \item No current or historical target input
      \end{itemize}
      \vspace{0.15cm}
      \small
      \begin{tabular}{lrr}
        \toprule
        Test measure & Mean baseline & Blend \\
        \midrule
        MAE  & 0.355 & \good{0.249} \\
        RMSE & 0.454 & \good{0.301} \\
        $R^2$ & $\approx 0$ & \good{0.558} \\
        \bottomrule
      \end{tabular}
    \column{0.41\textwidth}
      \textbf{Interpretation}
      \begin{itemize}
        \item RMSE improved by about \textbf{33.6\%} over the training-mean baseline.
        \item The linear and nonlinear models made complementary errors.
        \item An empirical 90\% interval achieved \textbf{89.7\%} test coverage.
        \item Performance degraded strongly outside the training envelope.
      \end{itemize}
  \end{columns}

  \vfill
  \alert{This estimates the current value; it is not yet a validated 10-minute early-warning forecast without target history.}
\end{frame}

\begin{frame}{9. The dashboard turns the notebooks into an investigation tool}
  \begin{columns}[T,onlytextwidth]
    \column{0.48\textwidth}
      \textbf{Available functions}
      \begin{itemize}
        \item Upload a CSV or use the default month
        \item Explore signals, missingness and correlations
        \item Detect candidate multivariate process-change windows
        \item Inspect which sensors drove each window
      \end{itemize}
    \column{0.48\textwidth}
      \textbf{Prediction functions}
      \begin{itemize}
        \item Target-anchored forecasts at 1, 5, 10 and 15 minutes
        \item Target-free virtual-sensor estimates
        \item Prediction intervals and drift flags
        \item Error metrics only when labels are available
      \end{itemize}
  \end{columns}

  \vspace{0.25cm}
  \begin{block}{How to read a detected window}
    The green score measures joint sensor movement, not target abnormality. A highlighted window is a candidate process transition; it must be compared with the later target response, prediction error and plant-event records.
  \end{block}
\end{frame}

\begin{frame}{10. Current conclusion and supervisor validation needs}
  \begin{columns}[T,onlytextwidth]
    \column{0.48\textwidth}
      \textbf{What the evidence supports}
      \begin{itemize}
        \item Target-anchored change forecasting beats persistence at all four horizons.
        \item The target-free virtual sensor is materially better than a constant baseline.
        \item Process-aware and temporal information add measurable value.
      \end{itemize}
    \column{0.48\textwidth}
      \textbf{What remains unproven}
      \begin{itemize}
        \item Generalization to another month, season or product grade
        \item Reliable operational warning during confirmed plant events
        \item Causal or controllable interpretation of important variables
      \end{itemize}
  \end{columns}

  \vspace{0.2cm}
  \begin{block}{Questions for process validation}
    Confirm units for acid, ground phosphate, fuel oil, steam, wash liquid and recyclage; define the recycled stream; confirm slurry mass/volume stream correspondence; identify startups, shutdowns and disturbances; confirm which variables are truly controllable.
  \end{block}

  \smallnote{Recommended next evidence: evaluate a later labeled month against the frozen pipelines and compare stable operation with confirmed transition events.}
\end{frame}

\end{document}
